% META: {"id":"1SPE-DERLOCAL-EV-B","chapitre":"1SPE-DERIVATION-LOCAL","type_objet":"evaluation","capacites_codes":["C1","C2","C3","C4","C5"],"duree_min":55,"points":20,"status":"generated","sympy_params":{"f_ex1":"x**2 - 6*x + 10","a_ex1":5,"f_ex2":"x**3","a_ex2":2,"f_ex3_a":3,"f_ex3_fa":4,"f_ex3_fpa":2,"f_ex4":"1/x","a_ex4":5}}

\section*{Évaluation B — Dérivation locale}

\textit{Durée : 55 minutes — 20 points — Calculatrice autorisée.}

\bigskip

%---------------------------------------------------------------
\textbf{Exercice 1} \hfill \textit{6 points — C1, C2}
%---------------------------------------------------------------

Soit $f$ la fonction définie sur $\mathbb{R}$ par $f(x) = x^2 - 6x + 10$.

\begin{enumerate}
  \item Calculer le taux de variation de $f$ entre $x = 2$ et $x = 5$. \hfill \textit{(1,5 pt)}
  \item Calculer $\dfrac{f(5+h)-f(5)}{h}$ pour $h \neq 0$. Simplifier. \hfill \textit{(2 pt)}
  \item En déduire le nombre dérivé $f'(5)$. \hfill \textit{(1 pt)}
  \item Interpréter $f'(5)$ graphiquement. \hfill \textit{(1,5 pt)}
\end{enumerate}

% BEGIN-VERIFY
% from sympy import symbols, expand, cancel, Rational
% x, h = symbols('x h')
% f = x**2 - 6*x + 10
% assert f.subs(x, 2) == 2
% assert f.subs(x, 5) == 5
% tv = Rational(5-2, 5-2)
% assert tv == 1
% f5 = f.subs(x, 5)
% assert f5 == 5
% f5h = expand(f.subs(x, 5+h))
% assert f5h == h**2 + 4*h + 5
% tvh = cancel((f5h - f5)/h)
% assert tvh == h + 4
% fp5 = tvh.subs(h, 0)
% assert fp5 == 4
% END-VERIFY

\bigskip

%---------------------------------------------------------------
\textbf{Exercice 2} \hfill \textit{5 points — C2, C4}
%---------------------------------------------------------------

Soit $g$ la fonction définie sur $\mathbb{R}$ par $g(x) = x^3$. On admet que $g'(a) = 3a^2$ pour tout $a \in \mathbb{R}$.

\begin{enumerate}
  \item Calculer $g'(2)$. \hfill \textit{(1 pt)}
  \item Écrire l'équation de la tangente $T_2$ à la courbe de $g$ au point d'abscisse $2$. \hfill \textit{(2 pt)}
  \item Vérifier que $T_2$ passe par le point $(0;-16)$. \hfill \textit{(1 pt)}
  \item Déterminer l'abscisse du point de la courbe où la tangente est horizontale. \hfill \textit{(1 pt)}
\end{enumerate}

% BEGIN-VERIFY
% from sympy import symbols, expand, solve
% x = symbols('x')
% g = x**3
% gp = 3*x**2
% assert gp.subs(x, 2) == 12
% assert g.subs(x, 2) == 8
% T2 = expand(12*(x-2) + 8)
% assert T2 == 12*x - 16
% assert T2.subs(x, 0) == -16
% assert solve(gp, x) == [0]
% END-VERIFY

\bigskip

%---------------------------------------------------------------
\textbf{Exercice 3} \hfill \textit{5 points — C3, C4, C5}
%---------------------------------------------------------------

On lit sur un graphique que la courbe de $f$ passe par le point $A(3;4)$ et que la tangente en $A$ a pour pente $2$.

\begin{enumerate}
  \item Donner les valeurs de $f(3)$ et $f'(3)$. \hfill \textit{(1 pt)}
  \item Écrire l'équation de la tangente en $A$. \hfill \textit{(1,5 pt)}
  \item Approcher $f(3{,}02)$ à l'aide de l'approximation linéaire. \hfill \textit{(1,5 pt)}
  \item Ce résultat est-il une valeur exacte ou approchée ? Justifier. \hfill \textit{(1 pt)}
\end{enumerate}

% BEGIN-VERIFY
% from sympy import symbols, expand, Rational
% x = symbols('x')
% T = expand(2*(x-3) + 4)
% assert T == 2*x - 2
% approx = 4 + 2*Rational(2,100)
% assert approx == Rational(404, 100)
% END-VERIFY

\bigskip

%---------------------------------------------------------------
\textbf{Exercice 4} \hfill \textit{4 points — C1, C5}
%---------------------------------------------------------------

Soit $h$ la fonction définie sur $]0;+\infty[$ par $h(x) = \dfrac{1}{x}$. On admet que $h'(a) = -\dfrac{1}{a^2}$ pour tout $a > 0$.

\begin{enumerate}
  \item Calculer le taux de variation de $h$ entre $x=4$ et $x=5$. \hfill \textit{(1,5 pt)}
  \item Approcher $h(5{,}02) = \dfrac{1}{5{,}02}$ à l'aide de l'approximation linéaire en $a=5$. \hfill \textit{(1,5 pt)}
  \item Comparer avec la valeur exacte $\dfrac{1}{5{,}02}$ (arrondir au millième). \hfill \textit{(1 pt)}
\end{enumerate}

% BEGIN-VERIFY
% from sympy import Rational
% h4 = Rational(1,4)
% h5 = Rational(1,5)
% tv = (h5-h4)/(5-4)
% assert tv == Rational(-1,20)
% hp5 = Rational(-1,25)
% approx = h5 + hp5*Rational(2,100)
% assert approx == Rational(249, 1250)
% exact = Rational(1, Rational(502,100))
% assert exact == Rational(50, 251)
% END-VERIFY
